\chapter{Zur kirchlichen Chronik}
\quellseite{8}

\section{America}

\subsection{Aus unserer Taubstummenmission}

Daß wir Taubstummenmissionare Zutritt haben zu einigen Staatsanstalten für Taubstumme und daselbst monatlich den Samen göttlichen Wortes ausstreuen unter eine Zuschauerschar von 200 bis 500 Schülern jeglichen Alters, dürfte wohl schon in weiteren Kreisen unserer Synode bekannt sein. Daß wir aber in diesen Schulen auch an Taubblinden missioniren, ist wohl weniger bekannt. Auf der Staatsschule von Wisconsin zu Delavan gibt es zwei solche arme, unglückliche Mädchen. Die eine, nun sechzehn Jahre alt, ist taubblind geboren. Die andere, jetzt siebzehn Jahre alt, wurde vor vier Jahren in Folge von Scharlachfieber taub und fast zur selben Zeit auch blind. Erstere fühlt sich weit glücklicher als ihre ältere Leidensgenossin, da sie nie erfahren hat, was es heißt, einen Laut zu hören, was es heißt, die Stimme einer liebevollen Mutter zu vernehmen. So hat sie auch in Folge ihrer frühen Blindheit keine Ahnung von der Schönheit ihrer Umgebung; sie weiß nicht, was es heißt, Gottes herrliche Natur in Augenschein zu nehmen. Letztere hingegen hat bis zu ihrem dreizehnten Lebensjahre ihre vollen fünf Sinne gehabt und vermißt in Folge dessen die verlorenen nun um so mehr. Wie unglücklich fühlt sich dies arme Menschenkind, wie niedergeschlagen, wie muthlos! Sie lag auf dem Krankenbett, als der Unterzeichnete das letzte Mal in der Staatsschule predigte. Auch ihre Leidensgenossin mußte zur selben Zeit einige Tage das Bett hüten. Beide hatten meine Anwesenheit in der Anstalt erfahren und ließen mich rufen. Wie freuten sie sich, als ich sie begrüßte, nämlich dadurch, daß ich in ihre Hände den Gruß „hineinbuchstabirte“. Wie freute sich besonders die Siebzehnjährige, wie herzhaft und liebevoll streichelte und drückte sie mir die Hand, als ich ihr von ihrem Heilande erzählte! „Du bist der erste“, sagte sie, „der mir nun wieder einmal etwas von Jesu sagt in den vier Jahren, seit ich taubblind bin. Früher habe ich in der Bibel gelesen und habe es gerne gethan. Meine Eltern waren norwegische Lutheraner; sie sind mir aber vorausgegangen in die Ewigkeit, und nun stehe ich fast ganz allein. Drei meiner Geschwister sind ebenfalls blind, aber nicht taub. Oft habe ich gewünscht, wieder einmal in der Bibel lesen zu können, aber nun ist mein Augenlicht gänzlich dahin. ... wie könnte“ — sie hielt inne, erschöpft von Erregung und unter Thränen mir die Hand drückend. Es läßt sich nicht recht beschreiben, wie begierig eine solche Person jeden Trost göttlichen Wortes aufnimmt, wie hell die Strahlen göttlichen Wortes in solch ein verkümmertes Menschenherz hineinscheinen. Immer mehr begehrte sie zu wissen von ihrem Jesus. Sie wünschte heimzugehen zu ihrem Heilande; doch ist es offenbar nicht Gottes Wille, sie schon jetzt abzurufen, denn sie steht ihrer Genesung entgegen. Nachdem ich ihr das Versprechen gegeben hatte, sie wieder aufzusuchen, nahm ich endlich von beiden Abschied, dem lieben Gott herzlich dankbar dafür, daß er uns Mittel und Wege gegeben hat, selbst den Taubblinden sein Wort verkündigen zu können. Dieser Fall zeigt aber auch wieder klar, daß Gott uns die Zeichensprache vornehmlich dazu gegeben hat, damit dadurch sein Reich gebaut werde und viele Seelen der armen Taubstummen selig werden. Er segne auch fernerhin unsere Taubstummenmission und verhelfe recht vielen dieser bedauernswerthen Menschen dorthin, „wo kein Leid, kein Geschrei, keine Schmerzen“ mehr sein werden, wo zum ewigen Lobe Gottes den Tauben das Gehör, den Blinden das Gesicht und den Stummen der Mund geöffnet sein wird.

\begin{unterschrift}
\unsicher{T. W.}
\end{unterschrift}

\emph{Die Logenreligion} von der Vaterschaft Gottes und der Bruderschaft der Menschen und von der Seligkeit durch Rechtthun wird jeden Tag vor den gesetzgebenden Versammlungen unsers Landes in den „Gebeten“ der Kapläne zum Ausdruck gebracht, theils ganz deutlich und mit klaren Worten, theils andeutungsweise. Es kann auch gar nicht anders sein. Denn wie soll doch in einem Senat oder in einem Abgeordnetenhause anders gebetet werden als in den Logen, da ja diese Körperschaften auch, gerade wie die Logen, aus Leuten allerlei Glaubens sich zusammensetzen? Da finden sich Katholiken und Protestanten, Juden und Mormonen, Gottesbekenner und Gottesleugner in bunter Mischung, und der Kaplan soll nun so „vorbeten“, daß es für Christen und Unchristen paßt. Damit stimmt, daß der jetzige Kaplan des Senates in Washington, Dr. C. C. Hale, ein Unitarier ist, der die Gottheit Christi leugnet und an keine Dreieinigkeit glaubt, und daß neulich sogar ein Jude, Rabbi D. Philipson, dort das Eröffnungsgebet sprach, während Hale die Sitzung mit dem Vater-Unser schloß. Die ganze Sache ist ein grober Unfug, eine Vermischung von Staat und Kirche, die gegen unsere Landesconstitution ist, und um so verwerflicher, als sie der heillosen Logenreligion, die schon genug durch die zahllosen geheimen Gesellschaften ausgebreitet wird, Eingang und ein gewisses Ansehen verleiht. Die meisten Mitglieder des Congresses werden wahrscheinlich ohnehin Logenbrüder sein. L. F.

\section{Ausland}

Der Einzelkelch bei der Abendmahlsfeier, der zuerst in den Vereinigten Staaten in Sectenkirchen eingeführt wurde, dann in England Eingang fand, kommt auch in Deutschland in Aufnahme. In Bremen ist in der Lamberti-Kirche der gemeinsame Kelch schon beseitigt, und eben jetzt wird aus Berlin berichtet, daß eine zahlreich besuchte Versammlung von Mitgliedern der St. Nikolai-Gemeinde nach einem Vortrage des Archidiaconus Seydel beschlossen hat, Einzelkelche beim Abendmahl einzuführen. Dieser Beschluß wird, wie zugleich mitgetheilt wird, voraussichtlich die Zustimmung des Kirchenvorstandes finden, um so mehr, als der Probst an St. Nikolai, Generalsuperintendent Dr. Faber, damit einverstanden ist. Seydel machte zugleich den Vorschlag, der ebenfalls angenommen wurde, daß die Gemeinde 30 bis 40 kleine Kelche anschaffen solle. Diese sollten vom Prediger bei der Feier mit dem gesegneten Weine gefüllt und den Communicanten gereicht werden. Nach der Benutzung sollten sie gereinigt und von neuem benutzt werden. Es stehe aber jedem frei, seinen eigenen Kelch mitzubringen. Die „Allgemeine Evangelisch-Lutherische Kirchenzeitung“ kann ihr Befremden darüber nicht verhehlen und findet es besonders auffallend, daß sich ein Generalsuperintendent dazu hergibt, eine „uralte, dem Christenvolke theure Sitte“ zu beseitigen. — Das
\quellseite{9}
heilige Abendmahl heißt und ist Communion, das heißt, Gemeinschaft. Alle, die am heiligen Abendmahle Theil nehmen, geben dadurch zu erkennen, daß sie auf das innigste zu einer Gemeinschaft verbunden sind. Bei der Einsetzung des heiligen Abendmahls nahm Christus den Kelch, gab ihn den Jüngern und sprach: „Trinket alle daraus“, und St. Paulus sagt: „Ein Brod ist’s; so sind wir viele Ein Leib, dieweil wir alle Eines Brodes theilhaftig sind.“ „Wir sind, durch Einen Geist, alle zu Einem Leibe getauft, wir seien Juden oder Griechen, Knechte oder Freie; und sind alle zu Einem Geist getränket“, 1 Cor. 10, 17. 12, 13. Diese Gemeinschaft und Verbindung wird äußerlich dargestellt durch das Trinken aus einem gemeinschaftlichen Kelch. Es ist auch bis jetzt noch kein einziger Fall von Ansteckung oder Krankheitsübertragung durch den gemeinschaftlichen Abendmahlskelch bewiesen.

Geheime Gesellschaften gibt es bekanntlich auch in der Heidenwelt, und besonders ist China davon heimgesucht. Die St. Louiser \emph{Review} theilt mit, daß es dort eine große Anzahl Logen gibt mit regelrechten Einführungsceremonien, Eiden, geheimen Zeichen und Paßwörtern. Zumeist tragen sie einen politischen Charakter, und es hat im 19. Jahrhundert kaum einen Aufstand gegen die Regierung oder einen Angriff gegen die Fremden gegeben, der nicht auf die chinesischen Geheimbündler zurückzuführen wäre. Der ganze Boxeraufstand vom Jahre 1900 hatte in ihrem Kreise seine Entstehung. Besonders sind auch die Freimaurer und Oddfellows unter den Chinesen vertreten.
